\section{相变的一般理论}

\subsection{序参量与 Landau 自由能}

连续相变附近，用序参量 $\eta$ 展开自由能密度（对称性允许的项）
\begin{equation}
  f(\eta,T)
  =f_0
  +a(T)\eta^2
  +b\eta^4
  +\cdots
  -h\eta,
\end{equation}
其中 $a(T)\approx a'(T-T_c)$，$b>0$。$h=0$ 时，$T>T_c$ 有唯一极小 $\eta=0$；$T<T_c$ 出现非零极小
\begin{equation}
  \eta_0=\pm\sqrt{-a/(2b)},
\end{equation}
对称性自发破缺。外场 $h$ 挑出方向并抹平奇异性。

\subsection{平均场与涨落}

平均场忽略空间涨落，给出经典临界指数。计入涨落后，引入 Ginzburg 判据：存在临界区，其中平均场失效，需用重整化群。下临界维数以下，涨落可完全破坏有序（如 1D Ising、$d\le2$ 的连续对称性与 Mermin--Wagner）。

\subsection{一级与二级}

自由能中若允许 $\eta^3$ 项或 $b$ 变号，可出现一级相变（潜热、序参量跳跃）。气液临界点与铁磁临界点在对称性提升后同属一类普适类。
